\documentclass[10pt,a4paper]{amsart}
\usepackage{amsmath,amssymb,amsthm,graphicx}
\RequirePackage{amssymb,amsmath,amsthm,bm}
\RequirePackage{thmtools}
\RequirePackage[capitalize]{cleveref}
\providecommand{\tightlist}{%
  \setlength{\itemsep}{0pt}\setlength{\parskip}{0pt}}

%%%%% varphi instead of phi
\renewcommand{\phi}{\varphi}

\mathchardef\mhyph="2D 
\providecommand{\A}{\mathcal{A}}
\providecommand{\C}{\mathcal{C}}
\providecommand{\L}{\mathcal{L}}
\providecommand{\B}{\mathcal{B}}
\providecommand{\M}{\mathcal{M}}
\renewcommand{\S}{\mathcal{S}}
\providecommand{\ol}[1]{\overline{#1}}
\providecommand{\mbf}[1]{\mathbf{#1}}
\providecommand{\mf}[1]{\mathfrak{#1}}
\providecommand{\mrm}[1]{\mathrm{#1}}
\providecommand{\mc}[1]{\mathcal{#1}}
\providecommand{\ra}{\rightarrow}
\providecommand{\Ra}{\Rightarrow}
\providecommand{\La}{\Leftarrow}
\providecommand{\lr}{\leftrightarrow}
\providecommand{\hf}{\mathbb{HF}}
\providecommand{\spec}{DgSp}
\providecommand{\bSigma}{\bm\Sigma}
\providecommand{\bDelta}{\bm\Delta}
\providecommand{\bPi}{\bm\Pi}
\providecommand{\Pinf}[1]{\Pi^{\mathrm{in}}_{#1}}
\providecommand{\Sinf}[1]{\Sigma^{\mathrm{in}}_{#1}}
\providecommand{\dSinf}[1]{d\mhyph\Sigma^{\mathrm{in}}_{#1}}
\providecommand{\Dinf}[1]{\Delta^{\mathrm{in}}_{#1}}
\providecommand{\bp}{\bar p}
\providecommand{\bq}{\bar q}
\providecommand{\ba}{\bar a}
\providecommand{\bb}{\bar b}
\providecommand{\bc}{\bar c}
\providecommand{\bd}{\bar d}
\providecommand{\bx}{\bar x}
\providecommand{\by}{\bar y}
\providecommand{\bz}{\bar z}
\providecommand{\LO}{\mathrm{LO}}
\declaretheorem[name={Theorem},numberwithin=section]{theorem}
\declaretheorem[name={Proposition}, sibling=theorem]{proposition}
\declaretheorem[name={Lemma}, sibling=theorem]{lemma}
\declaretheorem[name={Fact}, sibling=theorem]{fact}
\declaretheorem[name={Claim}, numberwithin=theorem]{claim}
\declaretheorem[name={Definition},numberwithin=section,sibling=theorem, style=definition]{definition}
\declaretheorem[name={Remark}, numbered=no,style=definition]{remark}
\declaretheorem[name={Corollary}, sibling=theorem]{corollary}
\declaretheorem[name={Question},sibling=theorem]{question}
\newcommand{\T}{{\mathcal{T}}}
\newcommand{\SR}{SR}
\newcommand{\pSR}{pSR}
\newcommand{\Q}{\mathbb{Q}}
\newcommand{\Z}{\mathbb{Z}}
\newcommand{\K}{\mathcal{K}}
\newcommand{\N}{\mathcal{N}}
\renewcommand{\L}{\mathcal{L}}
\providecommand{\Ginf}[1]{\Gamma^{\mathrm{in}}_{#1}}
\title{Equality of the Scott rank of a structure and of the Scott rank of its tree of tuples}
\author{David Gonzalez, Dino Rossegger}
\date{arxiv:2305.07126v1; CC BY 4.0}
\begin{document}
\maketitle
\noindent\textit{Editorial note.} The retained excerpt at source lines 45-53 ends in the middle of the article's sentence 'While the latter is not an ...'; the sentence continues at line 54, which points at the invariant table and is deliberately not extracted. The three rank definitions inside the excerpt are complete.\par\medskip
\noindent\textit{Editorial note.} The retained environments print with this package's own numbering, while the article numbers them within its sections; the original labels are preserved in the span metadata and no number is asserted here.\par\medskip
\noindent\textit{Editorial note.} The article's translation lemma at source lines 234-286 is a source-local support result and is delivered with its complete proof; it is not a newly counted target.\par\medskip
\noindent\textit{Editorial note.} No retained line contains a citation, a figure or an included file, so the package declares neither a bibliography nor an asset.\par\medskip
\noindent\textit{Editorial scope.} One original exact result: the equality of the Scott rank of a structure and of its tree of tuples, together with the article's own complete translation lemma on which the proof rests.\par\medskip
\noindent\textit{Allowed standard prerequisites (editorial).}\begin{itemize}
\item Infinitary formulas and back-and-forth relations. The complexity classes Pi\textasciicircum{}in\_alpha and Sigma\textasciicircum{}in\_alpha of L\_\{omega\_1,omega\} formulas, the back-and-forth relation =\_alpha, and Scott rank as the least ordinal whose Sigma\textasciicircum{}in formulas define all tuple orbits; all three ranks are defined in the retained context excerpt.
\item Labeled trees and the tree-of-tuples construction. The vocabulary of labeled trees, the finite atomic diagram coding used as a label, and the replication of every branch; the article's own translation lemma (source lines 234-286) is delivered with its complete proof and is not assumed.
\end{itemize}\medskip
In this article we give a comprehensive Scott analysis of the class of linear
orderings using three definitions of rank. Given a structure $\A$ the
\emph{(parameterless) Scott rank} $\SR(\A)$, is the least $\alpha$, such that
the automorphism orbits of all tuples in $\A$ are $\Sinf{\alpha}$ definable; the
\emph{parameterized Scott rank} $\pSR(\A)$, is the least ordinal $\alpha$ such
that there is a parameter $\bp\in A^{<\omega}$ such that $SR((\A,\bp))=\alpha$;
and the \emph{Scott sentence complexity} of $\A$, $SSC(\A)$ is the least of the
complexities $\Pinf{\alpha}$, $\Sinf{\alpha}$, $\dSinf{\alpha}$ such that $\A$
has a Scott sentence of one of these complexities. While the latter is not an

\subsection{The tree of tuples and Scott sentence complexity}
\begin{definition}\label{def:replabeledtreed}
   A \emph{replicated labeled tree} consists of a tree $(T,\succeq)$ with a parent function and
   labeling function $l:T\to \omega$ that satisfies that for every $\sigma\in T$
   with parent $\tau$, there exist infinitely many children $\tilde\sigma$ of
   $\tau$ such that $\T_\sigma\cong \T_{\tilde\sigma}$.
\end{definition}
\begin{definition}\label{def:tree-of-tuples}
   Given a structure $\A$ let $T_\A$ be the labeled tree consisting of all
   tuples from $\A$ ordered by inclusion where each tuple $\ba$ is labeled by a
   natural number coding its finite atomic diagram $D_\A(\ba)$ and the length of
   $\ba$. The \emph{tree of tuples $\T_\A$} is obtained by replicating every branch in
   $T_\A$ infinitely many times.
  \end{definition}

On the other hand, we start by showing that the Scott ranks of the tree of tuples of a structure
and the Scott rank of the structure are equal. Towards this we need a few
lemmas.

\begin{theorem}\label{thm:treeofstructuressr}
For any structure $\A$, $\SR(\A)=\SR(\T_\A)$. 
\end{theorem}

\begin{lemma}\label{lemma:translation} 
   Let $\A$ be a structure in vocabulary $\tau$.
   \begin{enumerate}
      \tightlist
      \item\label{eq:pushforward} Given a $\Pinf{\alpha}$ $\tau$-formula $\phi$, there is a
         $\Pinf{\alpha}$ formula $T_\phi$ in the vocabulary of labeled trees
         such that for any $\ba\in \A$ and every $\sigma\in \T_\A$ with
         $\sigma\succeq \ba$ 
         \[\A\models\phi(\ba)\iff \T_\A\models T_\phi(\sigma).\]
      \item\label{eq:pullback} Given a $\Pinf{\alpha}$ formula $\psi$ in the
         vocabulary of labeled trees, there is a $\Pinf{\alpha}$ $\tau$-formula
         $\psi^*$ such that for any tuple $\bar{\sigma}=(\ba_1,\dots,\ba_n)\in
         \T_\A$ \[\T_\A\models \psi(\bar\sigma) \iff
         \A\models\psi^*(\ba,\dots,\ba_n).\]
      \end{enumerate}
\end{lemma} 
\begin{proof}
   We first show \cref{eq:pushforward} starting with quantifier free formulas.
   Let $\phi$ be quantifier free with free variables $x_1,\dots x_n$ in
   disjunctive normal form. For each disjunct $\phi_i$ let $n_i$ be the index of
   the largest occuring variable and let $T_{\phi_i}$ be the formula expressing
   that the height of $x$ in the tree is at least $n_i$ and that
   $\phi_i\subseteq l(x)$. Then $T_\phi=\bigvee T_{\phi_i}$. The construction of
   $\T_\A$ guarantees that for $\sigma\succeq\ba\in \A$, $\T_\A\models
   T_\phi(\sigma)$ if and only if $\A\models \phi(\ba)$. Now, assume that we
   have a translation for $\phi(\bx,y)$ into $T_\phi(x)$. If $\psi=\exists y
   \phi(\bx,y)$, then let \[T_\psi= (\exists y\succeq x)\, T_\phi(y).\] We have
   that $\A\models \exists y\,\psi(\ba,y)$, if and only if there is $b\in \A$
   such that $\A\models\psi(\ba,b)$ if and only if there is $\sigma\succeq \ba
   \in \T_\A$ $\T_\A\models T_\psi(\sigma)$ if and only if $\T_\A\models
   T_\psi$. This proves the successor case. The limit case is immediate.

   We now prove \cref{eq:pullback} by defining a family of Turing-computable
   embeddings from the class of structures containing $\A$ with parameters to
   labeled trees with parameters. Each of these embeddings will in some sense
   extend the Friedman-Stanley embedding that acts without parameters. Fix a
   tuple $\bar\sigma=(\ba_1,\cdots,\ba_n)\in \T_\A$. Consider any structure $\B$
   in the class of $\A$ with parameters $\bp$ such that
   $|\bp|=|\ba_1\dots\ba_n|$. We now construct the embedding
   $\Psi_{\bar\sigma}$. Non-uniformly fix a partition
   $\bp=\bp_1,\dots\bp_n$ with $|\bp_i|=|\ba_i|$ and construct the tree of
   tuples of $\B$ as normal. To decide what elements become parameters in
   $\T_\B$, look for the lexicographical least elements $\tau_i$ coding the
   $\bp_i$ such that the finite substructure of $\T_\B$ described by $\tau$ is
   isomorphic to the one described by $\sigma$ in $\T_\A$. It is not difficult to check that this is a Turing computable
   embedding. Now, for any formula $\psi$ in the language of labeled trees, any
   $\B$ in the class of $\A$ and any $\bar\sigma=(\bb_0\dots\bb_n)\in \T_\B$ by
   the pull-back theorem for $\Psi_{\bar\sigma}$ there is $\psi^*$ of the same
   complexity as $\psi$ such that $\B\models \psi^*(\bb_0\dots\bb_n)$ if and
   only if $\T_\B\models \psi(\bar\sigma)$. As the image of $\Psi_{\bar\sigma}$
   contains the parameterized tree $(\T_\A,\bar\sigma)$, this proves
   \cref{eq:pullback}.
\end{proof}

\begin{proof}
Let $\A$ be a structure such that $SR(\T_\A)=\gamma$. Consider a tuple $\ba\in \A$. Let
$\psi\in\Sinf{\gamma}$ define the orbit of $\ba\in \T_\A$. By \cref{eq:pullback}
in \cref{lemma:translation} we have for all $\bb\in \A$ \[\A\models
\psi^*(\bb)\iff \T_\A\models \psi(\bb) \iff \bb\cong \ba.\] Therefore, $\psi^*$
defines the orbit of $\ba$ as desired and hence, $\SR(\A)\leq\SR(\T_\A)$. 

We now show $\SR(\T_\A)\leq\SR(\A)=\gamma$. Let $\bar\sigma=\ba_1\dots\ba_n$ be
a tuple in $\T_\A$. The automorphism orbit of each  $\ba_i$ is $\Sinf{\gamma}$
definable in $\A$, thus by \cref{eq:pushforward} in \cref{lemma:translation}
there is a $\Sinf{\gamma}$ formula, $\chi_i$,  true only of elements $\tau$ extending $\sigma$ 
that code elements automorphic to $\ba_i$ in $\A$. Now, let
$\psi$ describe the finite subtree containing precisely the elements of
$\bar\sigma$ as its leaves and let $\phi$ by the conjunction of $\psi$ together
with the $\chi_i$'s. Note that
because the labels encode the height of the labeled element in the tree, $\psi$
can be taken $\Sinf{1}$. Let $\bar \tau=\bb_1\dots\bb_n$ satisfy $\phi$. Then,
for every $\bb_i$, the subtrees $T_{\bb_i}$ rooted at $\bb_i$ are isomorphic to
the subtrees of $\ba_i$ and as the tree $\T_\A$ is replicated we can extend
these isomorphisms to an automorphism mapping $\bar\sigma\to \bar\tau$.
\end{proof}

\end{document}
